% TOP_PROBLEM: {"problemId": "problem.strong-compactness-versus-supercompactness-equiconsistency-problem", "canonicalTitle": "Strong Compactness versus Supercompactness Equiconsistency Problem", "releaseRank": 361, "primaryDomain": "logic_foundations_set_theory", "primaryDomainLabel": "Logic, foundations & set theory", "displayStatus": "open", "releaseStatus": "open"}
% !TeX program = lualatex
% ATTEMPT_STATUS: unresolved
\documentclass[11pt]{article}
\usepackage[margin=25mm]{geometry}
\usepackage{fontspec}
\setmainfont{Latin Modern Roman}
\usepackage{amsmath,amssymb,mathtools}
\usepackage{unicode-math}
\setmathfont{Latin Modern Math}
\usepackage{xurl,hyperref}
\hypersetup{colorlinks=true,urlcolor=blue}
\setcounter{secnumdepth}{0}
\setlength{\parindent}{0pt}
\setlength{\parskip}{5pt}
\setlength{\emergencystretch}{3em}
\begin{document}
\section*{\texorpdfstring{Strong Compactness versus Supercompactness Equiconsistency Problem}{Strong Compactness versus Supercompactness Equiconsistency Problem}}\label{361}
\subsection{Definitions and mathematical statement}
A cardinal $\kappa$ is strongly compact if every $\kappa$-complete filter extends to a $\kappa$-complete ultrafilter. It is supercompact if for every $\lambda\ge\kappa$ there is an elementary embedding $j:V\to M$ with critical point $\kappa$, $j(\kappa)>\lambda$, and $M^\lambda\subseteq M$, in the standard class-embedding formulation. Write $\operatorname{Con}(T)$ for the arithmetical assertion that the formal theory $T$ has no proof of contradiction. The requested reverse consistency implication is
\[
 \operatorname{Con}(\mathrm{ZFC}+\exists\text{ strongly compact cardinal})
 \ \Longrightarrow\
 \operatorname{Con}(\mathrm{ZFC}+\exists\text{ supercompact cardinal}).
\]
The forward large-cardinal implication supplies the other direction. The issue is equiconsistency, not whether every strongly compact cardinal in a given universe is supercompact. A precise strict-strength result also requires the usual specified metatheory for comparing consistency assertions.
\par\smallskip\textit{Formulation and concept references:} \href{https://arxiv.org/abs/1810.05058}{[S1]}.

\subsection{Short English statement}
Are the existence of a strongly compact cardinal and the existence of a supercompact cardinal equally strong consistency assumptions?
\subsection{Research attempt}
\paragraph{A necessary convention audit (27 September 2026).}
The usual large-cardinal definition of strongly compact includes that
$\kappa$ is uncountable. The printed definition does not say this.
Taken literally at $\kappa=\aleph_0$, a $\kappa$-complete filter is
just a filter closed under finite intersections, and every such proper
filter extends to an ultrafilter in ZFC by the ultrafilter lemma.
Thus the existence assertion on the left would already be a theorem
of ZFC under that literal reading.

This is a formulation issue, not a proposed solution of the conventional
large-cardinal problem. We preserve the inherited text and explicitly
use the standard uncountable-cardinal convention in the substantive
comparison below. Without that convention the displayed statement
instead compares the consistency of ZFC itself with a supercompact
extension. The two interpretations must not be silently interchanged.
We also understand filters to be proper; an improper filter cannot
extend to a proper ultrafilter.

Goldberg's primary paper [S1] proves structural implications under the
Ultrapower Axiom. Its abstract states the additional hypothesis and
exceptions. It supplies evidence for the consistency comparison, not
a theorem that ZFC alone establishes the requested reverse implication.
This attempt proves the easy consistency direction, analyzes the seed
and normality mechanism, and locates the precise missing uniformization
step. It does not assume the Ultrapower Axiom or its consistency strength.

\paragraph{The literal countable-cardinal problem can be isolated exactly.}
For completeness, extend a proper filter $\mathcal F$ on a set $I$
to a maximal proper filter by Zorn's lemma. A union of a chain of proper
filters is again a proper filter: every finite intersection is already
handled in one member of the chain, and the empty set belongs to none.
A maximal proper filter is an ultrafilter. If neither $A$ nor its
complement belonged to it, adjoining $A$ would either give a larger
proper filter or imply an old filter member disjoint from $A$; the
latter would put the complement of $A$ in the filter. Both alternatives
contradict the assumptions. This proves the ZFC extension theorem at
$\aleph_0$.

Hence, if countable $\kappa$ were allowed, the two theories
$\mathrm{ZFC}$ and ``ZFC plus a strongly compact cardinal'' as printed
would have the same theorems, up to that provable existence sentence.
Their arithmetical consistency statements would be provably equivalent.
This observation is separate from deciding an implication between their
truth values in an unspecified external universe.

It also has a proof-theoretic consequence under the usual hypotheses.
Let $T$ be the standard recursively axiomatized first-order formulation
of ZFC plus a supercompact cardinal. Such a cardinal is inaccessible,
so $T$ proves existence of a set model of ZFC and therefore
$\operatorname{Con}(\mathrm{ZFC})$. If ZFC proved the literal implication
$\operatorname{Con}(\mathrm{ZFC})\to\operatorname{Con}(T)$, then
$T$ would prove its own consistency. If $T$ is consistent, Gödel's
second incompleteness theorem rules that out. This is a conditional
nonprovability observation about the uncorrected reading, not a
metatheoretic disproof of the intended equiconsistency conjecture.

\paragraph{Consistency strength and implication inside one universe differ.}
Write $T_{\rm sc}$ and $T_{\rm sup}$ for the conventional uncountable
versions. Proving in ZFC that every supercompact is strongly compact
gives an effective translation of a contradiction proof in $T_{\rm sc}$
into one in $T_{\rm sup}$, by replacing the extra axiom with its theorem.
It follows, in a standard weak arithmetic metatheory adequate for coding
proofs, that
\[
 \operatorname{Con}(T_{\rm sup})\Longrightarrow
 \operatorname{Con}(T_{\rm sc}).
 \tag{361.1}
\]
The requested implication reverses (361.1). It does not assert that
an arbitrary strongly compact cardinal must be supercompact in the
same model. A forcing construction that changes the properties of one
cardinal while retaining a consistency upper bound is logically
compatible with equiconsistency. Conversely, a theorem identifying
cardinals after adding another axiom does not establish a consistency
comparison unless that extra axiom is itself obtained at the required
consistency strength.

The notation $j:V\to M$ is the standard shorthand for the embedding
or ultrafilter formulation of supercompactness. Consistency assertions
are taken for its usual first-order large-cardinal axiomatization,
not for an unrestricted second-order theory of all external class maps.
This prevents an accidental change of formal theories in applying
incompleteness or proof translations.

\paragraph{A full proof that supercompactness extends complete filters.}
Let $\kappa$ be supercompact, and let $\mathcal F$ be a proper
$\kappa$-complete filter on an arbitrary set $I$. Choose a cardinal
$\lambda$ at least the sizes of $I$ and $\mathcal F$, and an embedding
\[
 j:V\to M,\qquad \operatorname{crit}(j)=\kappa,
 \quad j(\kappa)>\lambda,\quad M^\lambda\subseteq M.
\]
We take $M$ transitive, as in the standard formulation. Closure puts
the set $j``\mathcal F$ in $M$: enumerate $\mathcal F$ in length at
most $\lambda$, take the sequence of its images, and use the closure
assumption. In $M$ this collection has size at most $\lambda<j(\kappa)$.

By elementarity $j(\mathcal F)$ is a proper $j(\kappa)$-complete
filter on $j(I)$. Every member of $j``\mathcal F$ belongs to it, so
its intersection belongs to it and is nonempty. Choose
\[
 a\in\bigcap_{A\in\mathcal F}j(A).
\]
Define a collection of subsets of $I$ by
\[
 U=\{A\subseteq I:a\in j(A)\}.
 \tag{361.2}
\]
Elementarity shows that $U$ is an ultrafilter: the seed lies in $j(I)$,
not in $j(\varnothing)$, intersections are respected, and exactly one
of $j(A),j(I\setminus A)$ contains it. The chosen intersection makes
$U$ extend $\mathcal F$.

For $\theta<\kappa$ and $A_\xi\in U$ for $\xi<\theta$, the critical
point condition fixes $\theta$ and all its indices. Therefore
\[
 j\left(\bigcap_{\xi<\theta}A_\xi\right)
      =\bigcap_{\xi<\theta}j(A_\xi),
\]
and the seed lies in this intersection. Thus $U$ is
$\kappa$-complete. This proves the full filter-extension implication,
with arbitrary $I$ and arbitrary filter size, and hence (361.1).
The dependence of the chosen embedding on $\mathcal F$ and $\lambda$
is allowed by supercompactness.

The proof identifies where the stronger hypothesis entered: closure
of $M$ placed the entire family of image sets inside $M$ with size
less than $j(\kappa)$. A single embedding without that closure would
not automatically permit this internal intersection argument.

\paragraph{Normal fine measures exhibit the extra seed information.}
For an uncountable regular $\kappa$ and $\lambda\ge\kappa$, put
\[
 P_\kappa(\lambda)=\{x\subseteq\lambda:|x|<\kappa\}.
\]
The cones $C_\alpha=\{x:\alpha\in x\}$ generate a proper
$\kappa$-complete filter when all intersections of fewer than $\kappa$
cones are included. Such an intersection asks that $x$ contain a
fixed set of size less than $\kappa$, which is possible. Regularity
ensures that the union of fewer than $\kappa$ such small sets is still
small. Extending this filter by strong compactness supplies a fine
$\kappa$-complete ultrafilter $U$; ``fine'' means $C_\alpha\in U$
for every $\alpha<\lambda$.

In contrast, a supercompact embedding supplies the specific seed
$j``\lambda\in P_{j(\kappa)}(j(\lambda))^M$ and the ultrafilter
\[
 U_j=\{A\subseteq P_\kappa(\lambda):j``\lambda\in j(A)\}.
 \tag{361.3}
\]
Closure ensures that this seed belongs to $M$, and it has size
$\lambda<j(\kappa)$ there. Fineness and $\kappa$-completeness follow
as in (361.2). In addition, it is normal: whenever a function $f$
on a $U_j$-large set satisfies $f(x)\in x$, then
$j(f)(j``\lambda)\in j``\lambda$, so it equals $j(\alpha)$ for
some $\alpha<\lambda$. The set of $x$ with $f(x)=\alpha$ then
belongs to $U_j$ by the definition of the seed measure.

Normality is this selection property for regressive functions, not
merely closure under fewer than $\kappa$ intersections. The filter
extension axiom promises an ultrafilter extending the cones; it does
not state that every such regressive choice becomes constant on a
large set. A proof of the reverse direction cannot import that
additional property by changing the adjective attached to the measure.

\paragraph{What the seed looks like in a general fine ultrapower.}
Suppose $U$ is a fine $\kappa$-complete ultrafilter on
$P_\kappa(\lambda)$, and let $j_U:V\to M_U$ be its well-founded
ultrapower. Countable completeness ensures well-foundedness; this is
the standard ultrapower fact used here. Write
\[
 s=[x\mapsto x]_U.
\]
Fineness implies $j_U``\lambda\subseteq s$, since each constant
$\alpha$ belongs to $x$ on a $U$-large cone. Also
$M_U\models |s|<j_U(\kappa)$, because every representative $x$
has size less than $\kappa$.

However these statements only give a small \emph{cover} of the image
of $\lambda$. They do not give equality. If $U$ is normal, equality
follows: every element of $s$ is represented, on a large set, by a
choice $f(x)\in x$, and normality makes it equal to a constant
$j_U(\alpha)$. Thus
\[
 U\text{ normal}\quad\Longrightarrow\quad s=j_U``\lambda.
 \tag{361.4}
\]
This proves the extra conclusion from its exact extra assumption.
Trying to prove (361.4) from fineness alone repeats the unresolved
selection step in a different notation.

The argument is diagnostic, not a claim that every fine complete
measure must be normal or that normalizing one isolated measure suffices
for full supercompactness. The consistency target requires a suitable
model with measures or embeddings at every relevant $\lambda$, and
all coherence and closure requirements must be justified there.

\paragraph{Why closure is much stronger than containing the image ordinals.}
Containing the set $j``\lambda$ is useful because it lets one index
sequences of images inside $M$. But closure $M^\lambda\subseteq M$
asserts that every external sequence of length $\lambda$ with values
in $M$ belongs to $M$, not merely the distinguished sequence of image
ordinals. Neither a cardinality estimate nor the existence of some
set $s$ covering $j``\lambda$ establishes this universal sequence
property by itself.

One standard route from normal fine measures to supercompactness proves
that sequence closure by working with representatives of the sequence
entries in the ultrapower. It uses normality to assemble the required
representatives along the small sets $x$. The elementary seed calculations
above do not perform this assembly for an arbitrary fine measure.
They locate the issue rather than announce that covering equals closure.

\paragraph{A direct check of the inaccessible-model consequence used earlier.}
We record the elementary critical-point facts to make the literal-reading
incompleteness diagnostic auditable. A supercompact critical point is
uncountable: an elementary embedding fixes every natural number and
$\omega$. It is regular. If a cofinal map $f:\theta\to\kappa$
existed with $\theta<\kappa$, then $j(f)$ would have the same domain
and the same values there, since all values are ordinals below the
critical point. Its supremum would be $\kappa$, contrary to elementarity,
which makes it cofinal in $j(\kappa)>\kappa$.

It is also a strong limit. For $\mu<\kappa$, every subset of $\mu$
is fixed by $j$, as is $\mu$. Transitivity then gives
$j(\mathcal P(\mu))=\mathcal P(\mu)$: all subsets in the source
are fixed and present in $M$, and $M$ has no additional external subsets.
If $2^\mu\ge\kappa$, choice would supply a surjection
$f:\mathcal P(\mu)\twoheadrightarrow\kappa$. Its image under $j$
has the same domain and agrees pointwise with $f$, so cannot be onto
$j(\kappa)$. This contradiction proves the strong-limit property.

Thus $\kappa$ is inaccessible. The standard rank-initial segment
$V_\kappa$ is a model of ZFC: power sets and the elementary set operations
stay below the limit rank; regularity and the bound $|V_\alpha|<\kappa$
for $\alpha<\kappa$ keep replacement images bounded in rank; choice
restricts to each set. This supplies the set-model consistency proof
used in the opening diagnostic. It is an implication available from a
supercompact cardinal, not a construction from bare consistency of ZFC.

\paragraph{Why the Ultrapower Axiom does not erase the consistency gap.}
Suppose a theorem has the form
\[
 \mathrm{ZFC}+\mathrm{UA}+\text{large-cardinal hypotheses}
       \vdash\text{a strong compactness/supercompactness comparison}.
\]
To infer the requested reverse consistency implication from it, one
would still need to obtain the joint theory including UA from the
consistency of $T_{\rm sc}$, with all exceptions and quantifiers in
the comparison accounted for. A conditional structural theorem alone
does not provide that consistency transfer. This is why the primary
paper is relevant evidence without being inserted as an unconditional
missing lemma.

Likewise a model in which some particular strongly compact cardinal
fails to be supercompact does not establish that the two existence
theories have different consistency strengths. Existence in another
model, possibly with a different cardinal playing the role, is what a
relative consistency construction would supply.

\paragraph{Outcome and first missing step.}
There is no meaningful finite-cardinality experiment that could verify
closure under all $\lambda$-sequences or decide consistency of these
infinite theories. Verification here consists of checking the quantified
filter and seed arguments, the direction of proof translations, and
the exact extra hypothesis in [S1]. Finite ultrafilter analogies are
not offered as evidence for a large-cardinal consistency theorem.

\textbf{Outcome: UNRESOLVED under the standard uncountable convention.}
The literal definition first needs that convention made explicit.
For the intended problem the easy consistency direction is proved,
and the fine-versus-normal seed distinction identifies a substantive
gap. No inner-model construction, forcing construction, or other
metatheoretic interpretation producing a supercompact-cardinal theory
from $\operatorname{Con}(T_{\rm sc})$ is established here.

\subsection{Sources}
\begin{itemize}
\item[S1]G. Goldberg, The Ultrapower Axiom and the equivalence between strong compactness and supercompactness, arXiv:1810.05058 (2018).
\url{https://arxiv.org/abs/1810.05058}.
\textit{Formulation source.}
\end{itemize}
\end{document}
